\subsection{Федеральный закон № 152-ФЗ «О персональных данных»}

Федеральный закон № 152-ФЗ «О персональных данных» устанавливает правовые основы обработки персональных данных и направлен на обеспечение защиты прав и свобод человека и гражданина при обработке его персональных данных, включая защиту прав на неприкосновенность частной жизни, личную и семейную тайну. Данный закон развивает положения Конституции Российской Федерации и конкретизирует требования к работе с данными, позволяющими идентифицировать субъекта.

Согласно статье 3 Федерального закона № 152-ФЗ персональные данные представляют собой любую информацию, относящуюся прямо или косвенно к определённому или определяемому физическому лицу (субъекту персональных данных), а обработка персональных данных включает любые действия с ними, в том числе сбор, запись, систематизацию, накопление, хранение, уточнение, использование и передачу. Оператором персональных данных является лицо, самостоятельно или совместно с другими лицами организующее и осуществляющее обработку данных, а также определяющее цели такой обработки.

В контексте разрабатываемого программного средства анализ audit-логов Kubernetes-кластера может затрагивать персональные данные, поскольку такие логи могут содержать идентификаторы пользователей, IP-адреса, временные метки, сведения о действиях в системе и другие атрибуты, позволяющие прямо или косвенно установить связь с конкретным субъектом. В соответствии со статьёй 5 закона обработка персональных данных должна осуществляться на законной и справедливой основе, ограничиваться достижением конкретных, заранее определённых и законных целей, а также не допускать избыточности обрабатываемых данных по отношению к заявленным целям.

Статья 6 Федерального закона № 152-ФЗ определяет условия обработки персональных данных, включая необходимость получения согласия субъекта либо наличие иных законных оснований для обработки. В случае использования системы обнаружения аномалий в рамках обеспечения информационной безопасности допускается обработка персональных данных без согласия субъекта, если это необходимо для достижения целей, предусмотренных законодательством, при условии соблюдения прав и законных интересов субъектов данных.

Особое значение имеет статья 19, устанавливающая обязанность оператора принимать необходимые правовые, организационные и технические меры для защиты персональных данных от неправомерного или случайного доступа, уничтожения, изменения, блокирования, копирования и распространения. К таким мерам относятся, в частности, разграничение прав доступа, использование средств защиты информации, а также контроль и аудит действий пользователей.

Таким образом, при разработке и внедрении системы обнаружения аномалий в Kubernetes-кластере необходимо учитывать требования Федерального закона № 152-ФЗ, включая минимизацию обрабатываемых данных, определение целей обработки, обеспечение законных оснований обработки и реализацию комплекса мер по защите персональных данных. Особое внимание должно уделяться возможности ограничению доступа к audit-логам, что позволяет снизить риски нарушения прав субъектов персональных данных.